% Reviewed translation: PDF pages 70--75 / printed pages 56--61.
% Begins after the completed Section 3 material on the shared PDF page 70.

\setcounter{section}{3}
\section{抽象变分问题的分析}\label{sec:abstract-variational-problem}

\MathBookSourcePage{70}
本节构造一个非常适合求解带约束的各类线性边值问题(如 Stokes 问题)的抽象框架.
我们提出若干处理约束的算法. 虽然这些算法是结合连续问题引入的,
但它们主要将在离散问题的求解中发挥作用.

\setcounter{subsection}{0}
\subsection{一般结果}\label{subsec:abstract-general-result}
\setcounter{equation}{0}
\setcounter{statement}{0}

\MathBookSourcePage{71}
设 $X$ 和 $M$ 是两个实 Hilbert 空间, 其范数分别为
$\|\cdot\|_X$ 和 $\|\cdot\|_M$. 设 $X'$ 和 $M'$ 为相应的对偶空间,
$\|\cdot\|_{X'}$ 和 $\|\cdot\|_{M'}$ 表示其对偶范数.
照例, 用 $\langle\cdot,\cdot\rangle$ 表示 $X$ 与 $X'$ 之间或
$M$ 与 $M'$ 之间的对偶配对.

引入两个连续双线性形式
\[
a(\cdot,\cdot):X\times X\longrightarrow\mathbb R,
\qquad
b(\cdot,\cdot):X\times M\longrightarrow\mathbb R,
\]
其范数为
\[
\|a\|=\sup_{\substack{u,v\in X\\u,v\ne0}}
\frac{a(u,v)}{\|u\|_X\|v\|_X},
\qquad
\|b\|=\sup_{\substack{v\in X,\ \mu\in M\\v\ne0,\ \mu\ne0}}
\frac{b(v,\mu)}{\|v\|_X\|\mu\|_M}.
\]
考虑如下变分问题, 称为问题 $(Q)$:

给定 $l\in X'$ 和 $\chi\in M'$, 求 $(u,\lambda)\in X\times M$, 使得
\begin{equation}\label{eq:abstract-q-first}
a(u,v)+b(v,\lambda)=\langle l,v\rangle
\qquad\forall v\in X,
\end{equation}
\begin{equation}\label{eq:abstract-q-second}
b(u,\mu)=\langle\chi,\mu\rangle
\qquad\forall\mu\in M.
\end{equation}

为了研究问题 $(Q)$, 还需要一些记号. 对双线性形式
$a(\cdot,\cdot)$ 和 $b(\cdot,\cdot)$, 分别定义线性算子
$A\in\mathcal L(X;X')$ 和 $B\in\mathcal L(X;M')$:
\begin{equation}\label{eq:abstract-operator-a}
\langle Au,v\rangle=a(u,v)
\qquad\forall u,v\in X,
\end{equation}
\begin{equation}\label{eq:abstract-operator-b}
\langle Bv,\mu\rangle=b(v,\mu)
\qquad\forall v\in X,\quad\forall\mu\in M.
\end{equation}
设 $B'\in\mathcal L(M;X')$ 为 $B$ 的对偶算子, 即
\begin{equation}\label{eq:abstract-dual-operator-b}
\langle B'\mu,v\rangle=\langle\mu,Bv\rangle=b(v,\mu)
\qquad\forall v\in X,\quad\forall\mu\in M.
\end{equation}
容易验证
\begin{equation}\label{eq:abstract-operator-norms}
\|A\|_{\mathcal L(X;X')}=\|a\|,
\qquad
\|B\|_{\mathcal L(X;M')}=\|b\|.
\end{equation}
利用这些算子, 问题 $(Q)$ 等价地写为: 求 $(u,\lambda)\in X\times M$, 使
\[
Au+B'\lambda=l\quad\text{in }X',
\qquad
Bu=\chi\quad\text{in }M'.
\]

现在定义线性算子 $\Phi\in\mathcal L(X\times M;X'\times M')$:
\[
\Phi(v,\mu)=(Av+B'\mu,Bv).
\]
如果 $\Phi$ 是从 $X\times M$ 到 $X'\times M'$ 的同构,
就称问题 $(Q)$ 适定. 我们的目的是导出问题 $(Q)$ 适定的充要条件.

\MathBookSourcePage{72}
置
\[
V=\Ker(B),
\]
更一般地, 对每个 $\chi\in M'$, 定义仿射流形
\[
V(\chi)=\{v\in X;\ Bv=\chi\}.
\]
等价地,
\begin{equation}\label{eq:abstract-affine-manifold}
\left\{
\begin{aligned}
V(\chi)&=\{v\in X;\ b(v,\mu)=\langle\chi,\mu\rangle
\quad\forall\mu\in M\},\\
V&=V(0).
\end{aligned}
\right.
\end{equation}
此外, 由算子 $B$ 的连续性可知 $V$ 是 $X$ 的闭子空间.

与问题 $(Q)$ 相联系, 定义如下问题, 称为问题 $(P)$: 求 $u\in V(\chi)$, 使
\begin{equation}\label{eq:abstract-problem-p}
a(u,v)=\langle l,v\rangle
\qquad\forall v\in V.
\end{equation}
显然, 如果 $(u,\lambda)\in X\times M$ 是问题 $(Q)$ 的解,
那么 $u\in V(\chi)$ 且 $u$ 是 
\eqref{eq:abstract-problem-p} 的解, 即问题 $(P)$ 的解.
我们希望找到适当条件, 保证这个命题的逆命题成立. 为此, 定义 $V$ 的极集
\[
V^0=\{g\in X';\ \langle g,v\rangle=0\quad\forall v\in V\}.
\]

\begin{Lemma}\label{lem:ch4-abstract-div}
下列三个性质等价:

\textup{(i)} 存在常数 $\beta>0$, 使
\begin{equation}\label{eq:abstract-inf-sup}
\inf_{\mu\in M}\sup_{v\in X}
\frac{b(v,\mu)}{\|v\|_X\|\mu\|_M}\ge\beta;
\end{equation}

\textup{(ii)} 算子 $B'$ 是从 $M$ 到 $V^0$ 的同构, 且
\begin{equation}\label{eq:abstract-bprime-lower-bound}
\|B'\mu\|_{X'}\ge\beta\|\mu\|_M
\qquad\forall\mu\in M;
\end{equation}

\textup{(iii)} 算子 $B$ 是从 $V^\perp$ 到 $M'$ 的同构, 且
\begin{equation}\label{eq:abstract-b-lower-bound}
\|Bv\|_{M'}\ge\beta\|v\|_X
\qquad\forall v\in V^\perp.
\end{equation}
\end{Lemma}
\begin{Proof}
1) 先证明性质 (i) 与 (ii) 等价. 由 
\eqref{eq:abstract-dual-operator-b}, 性质 (i) 意味着
\[
\|B'\mu\|_{X'}
=\sup_{\substack{v\in X\\v\ne0}}
\frac{\langle B'\mu,v\rangle}{\|v\|_X}
\ge\beta\|\mu\|_M
\qquad\forall\mu\in M,
\]
所以 
\eqref{eq:abstract-inf-sup} 等价于
\eqref{eq:abstract-bprime-lower-bound}. 因而 (ii) 蕴含 (i).
要证明 (i) 蕴含 (ii), 只需证明在条件
\eqref{eq:abstract-bprime-lower-bound} 下, $B'$ 是从 $M$ 到 $V^0$ 的同构.
显然, 由 
\eqref{eq:abstract-bprime-lower-bound}, $B'$ 是从 $M$ 到其值域
$\mathcal R(B')$ 的一一算子, 且其逆连续. 因此 $B'$ 是从 $M$ 到
$\mathcal R(B')$ 的同构, 从而 $\mathcal R(B')$ 是 $X'$ 的闭子空间.

\MathBookSourcePage{73}
于是只需证明
\[
\mathcal R(B')=V^0.
\]
为此应用 Banach 闭值域定理(见 Yosida [84]), 它断言
\[
\mathcal R(B')=(\Ker(B))^0=V^0.
\]
这证明了 1).

2) 现在证明性质 (ii) 与 (iii) 等价. 首先注意, $V^0$ 可与
$(V^\perp)'$ 等距等同. 事实上, 对 $v\in X$, 以 $v^\perp$ 表示
$v$ 在 $V^\perp$ 上的正交投影. 对每个 $g\in(V^\perp)'$, 定义
$\widetilde g\in X'$:
\[
\langle\widetilde g,v\rangle=\langle g,v^\perp\rangle
\qquad\forall v\in X.
\]
显然 $\widetilde g\in V^0$, 而且容易验证对应
$g\mapsto\widetilde g$ 是从 $(V^\perp)'$ 到 $V^0$ 的等距双射.
这使我们可以等同 $(V^\perp)'$ 与 $V^0$.

因此, $B$ 是从 $V^\perp$ 到 $M'$ 的同构且
\[
\|B^{-1}\|_{\mathcal L(M';V^\perp)}\le1/\beta,
\]
当且仅当 $B'$ 是从 $M$ 到 $(V^\perp)'=V^0$ 的同构且
\[
\|(B')^{-1}\|_{\mathcal L(V^0;M)}\le1/\beta.
\]
所以性质 (ii) 与 (iii) 等价.
\end{Proof}

条件 
\eqref{eq:abstract-inf-sup} 通常称为 ``inf-sup 条件'',
它由 Babuška [4] 和 Brezzi [13] 各自独立引入.
为陈述主要结果, 定义线性连续算子 $\pi\in\mathcal L(X';V')$:
\[
\langle\pi f,v\rangle=\langle f,v\rangle
\qquad\forall f\in X',\quad\forall v\in V.
\]
显然
\[
\|\pi f\|_{V'}\le\|f\|_{X'}.
\]

\setcounter{statement}{0}
\begin{Theorem}\label{thm:abstract-well-posedness}
问题 $(Q)$ 适定(即算子 $\Phi$ 是从 $X\times M$ 到 $X'\times M'$ 的同构),
当且仅当下列条件成立:

\textup{(i)} 算子 $\pi A$ 是从 $V$ 到 $V'$ 的同构;

\textup{(ii)} 双线性形式 $b(\cdot,\cdot)$ 满足 inf-sup 条件
\eqref{eq:abstract-inf-sup}.
\end{Theorem}
\begin{Proof}
1) 条件 (i) 和 (ii) 是充分的. 由 
\eqref{eq:abstract-inf-sup} 和 Lemma~\ref{lem:ch4-abstract-div},
存在唯一的 $u_0\in V^\perp$, 使
\[
Bu_0=\chi,
\qquad
\|u_0\|_X\le(1/\beta)\|\chi\|_{M'}.
\]

\MathBookSourcePage{74}
因此问题 $(P)$ 等价地写成: 求 $w=u-u_0\in V$, 使
\[
a(w,v)=\langle l,v\rangle-a(u_0,v)
\qquad\forall v\in V,
\]
或等价地满足
\[
\pi Aw=\pi(l-Au_0).
\]
因为 $\pi A$ 是从 $V$ 到 $V'$ 的同构, 问题 $(P)$ 有唯一解
$u=u_0+w\in V(\chi)$, 且
\[
\|w\|_X\le C_1\|\pi(l-Au_0)\|_{V'}
\le C_1\|l-Au_0\|_{X'},
\]
所以
\[
\|u\|_X\le C_2(\|l\|_{X'}+\|\chi\|_{M'}).
\]
现在 $l-Au\in V^0$. 因而由 Lemma~\ref{lem:ch4-abstract-div},
存在唯一的 $\lambda\in M$, 使
\[
B'\lambda=l-Au,
\]
并且
\[
\|\lambda\|_M\le(1/\beta)\|l-Au\|_{X'}
\le C_3(\|l\|_{X'}+\|\chi\|_{M'}).
\]
因此问题 $(Q)$ 有唯一解 $(u,\lambda)$, 且映射
$(l,\chi)\mapsto(u,\lambda)$ 从 $X'\times M'$ 到 $X\times M$ 连续.
这意味着 $\Phi$ 是从 $X\times M$ 到 $X'\times M'$ 的同构.

2) 条件 (i) 和 (ii) 是必要的. 假设 $\Phi$ 是从 $X\times M$ 到
$X'\times M'$ 的同构. 先证明 inf-sup 条件
\eqref{eq:abstract-inf-sup}. 令 $\chi\in M'$ 并置
$(u,\lambda)=\Phi^{-1}(0,\chi)$. 因为 $Bu=\chi$, 所以
$\mathcal R(B)=M'$. 因此 $B$ 是从 $V^\perp$ 到 $M'$ 的连续一一映射,
从而是从 $V^\perp$ 到 $M'$ 的同构. 由 Lemma~\ref{lem:ch4-abstract-div},
条件 
\eqref{eq:abstract-inf-sup} 成立.

接着证明 $\pi A$ 是从 $V$ 到 $V'$ 的同构. 先验证 $\pi A$ 在 $V$ 上单射.
设 $u\in V$ 满足 $\pi Au=0$, 则 $Au\in V^0$.
由于 inf-sup 条件成立, 由 Lemma~\ref{lem:ch4-abstract-div}, $B'$ 是从
$M$ 到 $V^0$ 的同构. 因而存在唯一的 $\lambda\in M$, 使
$B'\lambda=-Au$. 于是 $\Phi(u,\lambda)=(0,0)$, 从而 $u=0$.

现在验证 $\pi A$ 满射. 设 $g\in V'$. 由 Hahn--Banach Theorem,
至少存在一个 $l\in X'$, 使 $g=\pi l$. 置
$(u,\lambda)=\Phi^{-1}(l,0)$. 显然 $u\in V$, 且
\[
Au+B'\lambda=l.
\]
对所有 $v\in V$,
\[
\langle\pi B'\lambda,v\rangle
=\langle B'\lambda,v\rangle
=\langle\lambda,Bv\rangle=0,
\]
所以 $\pi B'\lambda=0$, 从而

\MathBookSourcePage{75}
\[
\pi Au=\pi l=g.
\]
因此 $\pi A$ 是从 $V$ 到 $V'$ 的一一线性连续映射, 从而是同构.
\end{Proof}

\setcounter{statement}{0}
\begin{Remark}\label{rem:abstract-problem-p-solvability}
在 Theorem~\ref{thm:abstract-well-posedness} 证明的第一部分中,
我们使用了如下事实: 只要 $\pi A$ 是从 $V$ 到 $V'$ 的同构且仿射流形
$V(\chi)$ 非空, 问题 $(P)$ 就有唯一解. Inf-sup 条件
\eqref{eq:abstract-inf-sup} 确实蕴含 $V(\chi)$ 非空.
\end{Remark}

\setcounter{statement}{0}
\begin{Corollary}\label{cor:abstract-v-elliptic-well-posedness}
假设双线性形式 $a(\cdot,\cdot)$ 是 $V$-椭圆的, 即存在常数 $\alpha>0$, 使
\begin{equation}\label{eq:abstract-v-ellipticity}
a(v,v)\ge\alpha\|v\|_X^2
\qquad\forall v\in V.
\end{equation}
则问题 $(Q)$ 适定, 当且仅当双线性形式 $b(\cdot,\cdot)$ 满足 inf-sup 条件
\eqref{eq:abstract-inf-sup}.
\end{Corollary}
\begin{Proof}
设 $l\in V'$. 因为 $a(\cdot,\cdot)$ 是 $V$-椭圆的,
可以应用 Lax \& Milgram Theorem~\ref{thm:lax-milgram}:
存在唯一的 $u\in V$, 使
\[
a(u,v)=\langle l,v\rangle
\qquad\forall v\in V,
\]
或等价地
\[
\pi Au=l.
\]
此外, 映射 $l\mapsto u$ 从 $V'$ 到 $V$ 连续. 因此 $\pi A$ 是从
$V$ 到 $V'$ 的同构, 结论由 Theorem~\ref{thm:abstract-well-posedness} 得到.
\end{Proof}

\setcounter{statement}{1}
\begin{Remark}\label{rem:abstract-reflexive-banach-extension}
上述分析容易推广到 $X$ 和 $M$ 为自反 Banach 空间的情形.
空间 $V$ 不变, 但 $V^\perp$ 要换成商空间 $X/V$.
本小节的结果和证明只需极少修改即可适用.
\end{Remark}

\setcounter{subsection}{1}
\subsection{鞍点方法}\label{subsec:abstract-saddle-point}

在适当假设下, 问题 $(P)$ 和 $(Q)$ 可以表述为优化问题.
除上一小节的记号外, 再引入两个二次泛函
$J:X\to\mathbb R$ 和 $\mathcal L:X\times M\to\mathbb R$:
\begin{equation}\label{eq:abstract-functional-j}
J(v)=(1/2)a(v,v)-\langle l,v\rangle,
\end{equation}
以及
\begin{equation}\label{eq:abstract-lagrangian}
\mathcal L(v,\mu)=J(v)+b(v,\mu)-\langle\chi,\mu\rangle.
\end{equation}
